% This is samplepaper.tex, a sample chapter demonstrating the
% LLNCS macro package for Springer Computer Science proceedings;
% Version 2.21 of 2022/01/12
%
\documentclass{llncs}
%
\usepackage[T1]{fontenc}
% T1 fonts will be used to generate the final print and online PDFs,
% so please use T1 fonts in your manuscript whenever possible.
% Other font encondings may result in incorrect characters.
%
\usepackage{graphicx}
% Used for displaying a sample figure. If possible, figure files should
% be included in EPS format.
%
% If you use the hyperref package, please uncomment the following two lines
% to display URLs in blue roman font according to Springer's eBook style:
%\usepackage{color}
%\renewcommand\UrlFont{\color{blue}\rmfamily}
%\urlstyle{rm}
%

\usepackage[style=apa]{biblatex} % <- AGREGAR ESTA LÍNEA
\addbibresource{thebiblio.bib} % <- AGREGAR ESTA LÍNEA
\begin{document}
%
% ENGLISH VERSION
\title{Speckle Noise Removal in SAR Images using Autoencoders}

%\titlerunning{Abbreviated paper title}
% If the paper title is too long for the running head, you can set
% an abbreviated paper title here
%
\author{Bianca Balzarini\inst{1}\orcidID{0000-1111-2222-3333} \and
Juliana Gambini\inst{2,3,1}\orcidID{0000-0002-4534-1402}}
%
\authorrunning{B. Balzarini et al.}
% First names are abbreviated in the running head.
% If there are more than two authors, 'et al.' is used.
%
\institute{Instituto Tecnológico de Buenos Aires, Bueno Aires, Argentina
\email{bbalzarini@itba.edu.ar}\\
 \and
LIDEC, Universidad Nacional de Hurlingham, Pcia. de Buenos Aires, Argentina\\
\email{juliana.gambini@unahur.edu.ar} \\
\and
CPSI, Universidad Tecnológica Nacional, Regional Buenos Aires}
%
\maketitle              % typeset the header of the contribution
%
\begin{abstract}
Synthetic Aperture Radar (SAR) images are fundamental for Earth observation due to their ability to acquire data independently of weather and lighting conditions. However, these images are affected by speckle noise, a multiplicative granular pattern that degrades their quality and makes them difficult to automatically interpret.

This work presents an approach for speckle removal using convolutional autoencoders. The main innovation lies in the use of the $\mathcal{G}_I^0$ distribution  to synthetically generate pairs of noisy and clean images for supervised training, adequately modeling the statistical characteristics of real speckle noise.
In addition, we developed a flexible file-based configuration system that facilitates experimentation with different network architectures.

Experimental outcomes reveal that deeper network architectures achieve better results with no evidence of overfitting.

Models trained on synthetic data showed good performance on synthetic test images, but limitations in generalizing to actual SAR images. To address this problem, we explored training with pairs of real images obtained through time averaging, achieving significant improvements and performance comparable to other known models.

% keywords in lowercase
\keywords{Aperture Synthetic Radar \and Speckle noise \and Autoencoders}
\end{abstract}

% SPANISH VERSION
% Clear the previous title info and set Spanish version

% Spanish title and abstract
\begin{center}
\vspace{1cm}
{\Large \bfseries\boldmath Eliminación de Ruido \emph{Speckle} en Imágenes SAR utilizando Autoencoders}

\vspace{0.5cm}
\end{center}

\begin{abstract}
Las imágenes de Radar de Apertura Sintética (SAR: Synthetic Aperture Radar) son fundamentales para la observación
de la Tierra debido a su capacidad de adquisición independiente de las condiciones
climáticas y de iluminación. Sin embargo, estas imágenes se ven afectadas por el ruido \emph{speckle},
un patrón granular multiplicativo que degrada su calidad, y dificulta su interpretación.
Este trabajo presenta un enfoque para la eliminación de ruido \emph{speckle} mediante
el uso de autoencoders convolucionales. La principal innovación radica en el uso de la distribución
$\mathcal{G}_I^0$ para generar sintéticamente pares de imágenes ruidosas y limpias destinadas
al entrenamiento supervisado, modelando adecuadamente las características estadísticas del ruido
\emph{speckle} real. Complementariamente, desarrollamos un sistema flexible basado en archivos
de configuración que facilita la experimentación con diferentes arquitecturas de redes.
Los experimentos revelan que las arquitecturas más profundas obtienen consistentemente
mejores resultados sin evidencia de overfitting.

 Los modelos entrenados con datos sintéticos
mostraron buen desempeño en imágenes de prueba sintéticas, pero limitaciones en la generalización
a imágenes SAR reales. Para abordar este problema, exploramos el entrenamiento con pares de
imágenes reales obtenidas mediante promediado temporal, logrando mejoras significativas
y desempeño comparable a otros ya conocidos.
%Este trabajo establece las bases para el uso de autoencoders en el filtrado de ruido \emph{speckle} y
%demuestra tanto el potencial como las limitaciones de los datos sintéticos generados mediante
%modelos estadísticos.


% palabras clave en minúsculas
\keywords{ Radar de Apertura Sintética (SAR)
\and  ruido speckle\and Autoencoders}
\end{abstract}


\vspace{0.5cm}
\setcounter{page}{1}

%
%\section{First Section}
%\subsection{A Subsection Sample}
%Please note that the first paragraph of a section or subsection is
%not indented. The first paragraph that follows a table, figure,
%equation etc. does not need an indent, either.
%
%Subsequent paragraphs, however, are indented.
%
%\subsubsection{Sample Heading (Third Level)} Only two levels of
%headings should be numbered. Lower level headings remain unnumbered;
%they are formatted as run-in headings.
%
%\paragraph{Sample Heading (Fourth Level)}
%The contribution should contain no more than four levels of
%headings. Table~\ref{tab1} gives a summary of all heading levels.
%
%\begin{table}
%\caption{Table captions should be placed above the
%tables.}\label{tab1}
%\begin{tabular}{|l|l|l|}
%\hline
%Heading level &  Example & Font size and style\\
%\hline
%Title (centered) &  {\Large\bfseries Lecture Notes} & 14 point, bold\\
%1st-level heading &  {\large\bfseries 1 Introduction} & 12 point, bold\\
%2nd-level heading & {\bfseries 2.1 Printing Area} & 10 point, bold\\
%3rd-level heading & {\bfseries Run-in Heading in Bold.} Text follows & 10 point, bold\\
%4th-level heading & {\itshape Lowest Level Heading.} Text follows & 10 point, italic\\
%\hline
%\end{tabular}
%\end{table}
%
%
%\noindent Displayed equations are centered and set on a separate
%line.
%\begin{equation}
%x + y = z
%\end{equation}
%Please try to avoid rasterized images for line-art diagrams and
%schemas. Whenever possible, use vector graphics instead (see
%Fig.~\ref{fig1}).
%
%\begin{figure}
%\includegraphics[width=\textwidth]{fig1.eps}
%\caption{A figure caption is always placed below the illustration.
%Please note that short captions are centered, while long ones are
%justified by the macro package automatically.} \label{fig1}
%\end{figure}
%
%\begin{theorem}
%This is a sample theorem. The run-in heading is set in bold, while
%the following text appears in italics. Definitions, lemmas,
%propositions, and corollaries are styled the same way.
%\end{theorem}
%%
%% the environments 'definition', 'lemma', 'proposition', 'corollary',
%% 'remark', and 'example' are defined in the LLNCS documentclass as well.
%%
%\begin{proof}
%Proofs, examples, and remarks have the initial word in italics,
%while the following text appears in normal font.
%\end{proof}
%For citations of references, we use the APA 7 convention.  The following bibliography provides
%a sample reference list with entries for journal
%articles~\parencite{ref_article1}, an LNCS chapter~\parencite{ref_lncs1}, a
%book~\parencite{ref_book1}, proceedings without editors~\parencite{ref_proc1},
%and a homepage~\parencite{ref_url1}. Multiple citations are grouped
%\parencite{ref_article1,ref_lncs1,ref_book1},
%\parencite{ref_article1,ref_book1,ref_proc1,ref_url1}.
%
%\begin{credits}
%\subsubsection{\ackname} A bold run-in heading in small font size at the end of the paper is
%used for general acknowledgments, for example: This study was funded
%by X (grant number Y).
%
%\subsubsection{\discintname}
%It is now necessary to declare any competing interests or to specifically
%state that the authors have no competing interests. Please place the
%statement with a bold run-in heading in small font size beneath the
%(optional) acknowledgments,
%for example: The authors have no competing interests to declare that are
%relevant to the content of this article. Or: Author A has received research
%grants from Company W. Author B has received a speaker honorarium from
%Company X and owns stock in Company Y. Author C is a member of committee Z.
%\end{credits}
%%
%% ---- Bibliography ----
%%

%\printbibliography

\end{document}